\documentclass[11pt]{article}
\usepackage[utf8]{inputenc}
\usepackage[T1]{fontenc}
\usepackage{ebgaramond}
\usepackage[margin=1in]{geometry}
\usepackage{setspace}
\usepackage{titlesec}
\usepackage{hyperref}

\onehalfspacing

\titleformat{\section}[block]{\centering\large\scshape}{}{0em}{}
\titlespacing*{\section}{0pt}{2em}{1.5em}

\title{\textbf{Hemorrhagic}}
\author{Alex Towell}
\date{}

\begin{document}
\maketitle
\thispagestyle{empty}

\section*{The Lab That Was}

The centrifuges were under plastic sheeting. Amara passed them every morning on her way to the sequencing bench, and every morning she did not look at them directly, the way you learn not to look at the empty chair after a funeral.

Six weeks since the moratorium. The BSL-3 lab she had spent five years rebuilding---grant applications, contractor arguments, three trips to Geneva for equipment certifications---was dark. The negative-pressure seals still held. The HEPA filters were still rated. The biosafety cabinets stood ready behind their plastic shrouds, like surgical instruments laid out for an operation that had been cancelled.

Her two postdocs had been reassigned. David was in Accra now, running antibody panels for a measles surveillance project. Fatoumata had gone back to Dakar. Neither had argued. They understood the math. Everyone understood the math.

She had read the analysis herself---the one from the American artificial intelligence system, the one the world had been arguing about since it started issuing policy recommendations three months ago. Twenty-three percent probability of a lab-origin pandemic within a decade. Expected casualties in the millions. Amara had spent twenty years in biosafety. She had seen the mislabeled vials, the containment breaches noted in logbooks and never reported upward, the postdoc at the Pasteur Institute who had pipetted a live sample with a torn glove and not realized until he was halfway through decontamination. Twenty-three percent felt low, if anything.

The moratorium was defensible. She would say this to anyone who asked, and mean it.

But every morning she passed the dark lab, and every morning the plastic sheeting looked more permanent.

\begin{center}\rule{0.3\linewidth}{0.4pt}\end{center}

The farmer arrived at Redemption Hospital on a Tuesday. Forty-one years old, from Bong County, presenting with high fever and pain behind the eyes. No known exposure to Ebola or Marburg. No travel history beyond the usual route between his farm and the Gbarnga market.

Amara ran the panel. Negative for Ebola, Marburg, Lassa, dengue, chikungunya. She sequenced the sample herself because the technician was on leave and because she wanted to see it.

Novel paramyxovirus. Hemorrhagic features. She stared at the receptor binding profile for a long time. Then she enlarged the spike protein model and stared at that.

Respiratory tropism. High binding affinity. The virus was shaped to enter human lung tissue the way a key enters a lock --- not forcing its way in but fitting, as though it had been designed for this, which of course it had, by four billion years of evolution optimizing for exactly this moment, this receptor, this host.

Her mind did what it always did. It ran ahead.

Respiratory spread meant airborne transmission. Airborne transmission in Monrovia meant the buses, the markets, the churches where people stood shoulder to shoulder for two hours every Sunday. It meant the dense neighborhoods along the Mesurado where families of eight shared two rooms and the drainage ran open between the houses. It meant that by the time she identified the second case, there would already be a tenth.

And the incubation period --- she pulled up the clinical timeline from the farmer's admission records, cross-referenced with the blood panels --- looked long. Days, not hours. Days during which a person felt nothing worse than a mild headache, maybe a low fever they attributed to weather or overwork, while the virus replicated in their respiratory epithelium and shed with every breath and every cough and every word spoken to a neighbor over a shared plate of rice.

She set down the stylus. The sequence glowed on the screen. She could see, with the clarity of someone who had published forty-seven papers on viral emergence, exactly what was about to happen.

She reached for protocols that no longer existed.

The attenuated-virus pathway --- the fast one, the approach that had produced the Ebola VSV-ZEBOV vaccine in weeks instead of months --- required gain-of-function techniques. You took the wild virus, engineered it to replicate without causing disease, and used the weakened version to train the immune system. Standard practice. The equipment was three doors down the hall, under plastic.

She could see the equipment from where she sat. Through the lab's interior window. The centrifuges, the biosafety cabinet, the PCR thermal cycler she had calibrated herself. Dark shapes under translucent sheeting, like furniture in a house whose owners had gone and were not coming back.

She filed the emergency exception request. International treaty framework. Legal review required. Estimated processing time: six to eight weeks.

The virus did not observe treaty frameworks.

She called the WHO regional office in Brazzaville. Reported the novel isolate. Began conventional containment protocols. Contact tracing. Isolation of the index case. Notification of district health officers in Bong County. The procedures were automatic. She had done this before. She could do this in her sleep, and lately she nearly had been.

One case. Manageable.

She did not believe this, but she wrote it in the log.

\begin{center}\rule{0.3\linewidth}{0.4pt}\end{center}

Samuel brought lunch at one o'clock, the way he always did during outbreak responses. Rice and cassava leaf in a thermos, still hot. He taught secondary school mathematics three blocks from the hospital, and his free period aligned with the hour Amara was most likely to forget that eating was something her body required.

``You have your face on,'' he said, setting the thermos on her desk beside the sequencing printouts.

``I always have my face on.''

``Your outbreak face. The one where you look at people but you're counting something.''

She opened the thermos. The smell of palm oil and scotch bonnet filled the small office, and for a moment the room was not a room adjacent to a virology lab in a country where something terrible was starting to happen. It was just lunch. ``Novel paramyxovirus. One case. Probably nothing.''

Samuel sat in the chair he had sat in during the Ebola response, when ``probably nothing'' had become two thousand dead in Montserrado County alone. He did not say this. He watched her eat.

``Kofi needs help with quadratic equations,'' he said. ``He texted you.''

She had not checked her phone since the sequencing results. She picked it up. Three messages from Kofi: a photograph of his homework, a string of question marks, and then \emph{never mind papa explained it}. Below that, a photograph from Kwame's teacher: her youngest had drawn a picture titled ``Mama Fighting Germs.'' Stick figure in a white coat, holding something that could be a sword or a pipette, facing a crowd of green circles with angry faces.

``He drew teeth on the germs,'' she said.

``He draws teeth on everything. Last week he drew teeth on the sun.''

She ate. Samuel talked about the school football match, about the leak in the kitchen ceiling that the landlord had promised to fix for the third time, about Ama's request for a mobile phone because every other girl in her class apparently had one. She let his voice fill the room. This was what he did --- not because he didn't sense the gravity of what was on her screen, but because he understood that the gravity would still be there after lunch, and the rice would not.

After he left, Amara washed the thermos in the lab sink and set it beside the centrifuge covers. She checked the farmer's vitals. Stable, for now. Temperature 39.2, coming down with antipyretics.

Then she went back to the sequence on her screen and began designing a conventional vaccine candidate. The slow way. The way that worked, eventually, if you had enough time.

She did not have enough time. She knew this already. But she began.

\section*{The Parking Lot}

Three became twelve. Twelve became forty. Forty became a number that changed every time she refreshed the database.

Bong County. Kakata. Then the Monrovia suburbs --- Paynesville first, then Sinkor, then the dense neighborhoods along the Mesurado. The virus moved the way she had known it would move, following the bus routes and the market crowds and the Sunday church services, spreading through handshakes and shared meals and the thousand small intimacies of a city where people lived close because they had no choice.

The incubation period was the thing that made this different from Ebola. Ebola announced itself fast --- you were visibly ill within days, too sick to move far, too obviously dying to hide. This virus gave you a week. A week of mild symptoms that looked like malaria or a bad cold, during which you rode the bus to Waterside Market, sat in the pew next to your grandmother, held your neighbor's baby while she cooked. By the time the hemorrhagic symptoms appeared --- the blood in the urine, the petechiae spreading up the forearms, the cough that left iron on the tongue --- you had been breathing on everyone you loved for seven days.

Amara watched it happen with the clarity of someone who had seen it before. Not this virus, but this pattern. The same curve, the same lag between what was happening and what people believed was happening. The same moment --- there was always a moment --- when the exponential became visible to everyone, not just the epidemiologists, and by then it was too late to do anything except manage the dying.

That moment came on the fourth day. Redemption Hospital's isolation ward had twelve beds. Amara had designed the ward herself after 2014, argued with the contractors about HEPA specifications, tested the negative-pressure seals by holding a tissue to the door gaps and watching it pull inward. Twelve beds, properly ventilated, properly sealed, the best isolation capacity in the country.

By the fifth day, she needed thirty.

The parking lot became a ward. Not by decision --- no one decided --- but by accumulation. A patient arrived and there was no bed. A nurse set up an IV stand next to the ambulance bay. Someone found plastic sheeting. Someone else carried a mattress through the double doors. Within two days, the car park held more patients than the building. Plastic sheeting over concrete barriers. IV poles improvised from broom handles and electrical conduit. Extension cords snaking from the generator to a cluster of monitors that beeped under a tarpaulin while rain pooled on the sagging plastic above them.

The generator ran twenty hours a day. When it failed --- and it failed twice in the first week, once because the fuel delivery was late, once because the voltage regulator burned out --- the ventilators stopped.

The first time, the silence was the worst part. The constant mechanical breathing ceased and in its place was the sound of people trying to breathe on their own. Some of them could. The woman in her sixties from Sinkor could not. Amara heard the change in her breathing --- the wet, effortful sound of lungs filling with what lungs should not fill with --- and ran to the bedside, but the ventilator was off and the backup generator was thirty seconds from engaging and thirty seconds was enough. The woman's name was Mary Kollie. She had been stable.

The second time, Amara was already standing beside the generator when it failed, because she had learned. She switched to the backup herself, hands slippery with sweat, and the ventilators resumed in eleven seconds. No one died. But the boy from Paynesville --- sixteen, had been improving, his mother had brought him groundnuts that morning --- desaturated during those eleven seconds and never recovered. He died the next day, his oxygen levels chasing a threshold his damaged lungs could no longer reach.

Two patients. Not killed by the virus. Killed by a fuel truck delayed by traffic on Tubman Boulevard. Killed by a voltage regulator that had been in service since the Ebola rebuild and should have been replaced two years ago. Killed by a supply chain that ran from a manufacturer in Shenzhen through a distributor in Accra through a procurement office in Monrovia that had been operating on emergency budgets since 2014 and had never fully restocked.

Amara noted their deaths in the log.

Then she walked down the hall, past the dark BSL-3 lab with its covered centrifuges, and locked herself in the staff bathroom and pressed her forehead against the tile wall. The tile was cool. She counted to thirty. She did not cry because she did not have time to recover from crying, and she had four hours left on her shift and seven new admissions pending.

She went back to the ward.

\begin{center}\rule{0.3\linewidth}{0.4pt}\end{center}

Rebecca Foster arrived on the sixth day with an MSF advance team of four. Young --- thirty-one, Amara would learn later --- with the particular exhaustion of someone who had just finished one deployment and volunteered for the next without going home.

They worked well together. You learned quickly whether someone could work in a hot zone or whether they would hesitate at the wrong moment, and Rebecca did not hesitate. She had steady hands and a habit of narrating her procedures under her breath --- ``gown, gloves, first pair, tape, second pair, tape, face shield, check seal'' --- a liturgy of self-preservation that Amara recognized because she used the same one. The ritual that kept you alive when your hands wanted to rush and your body wanted to be anywhere else.

On the eighth day, Rebecca found Amara in the staff room, staring at the laptop. The emergency exception request. Status: \emph{Under review --- Legal Advisory Panel, International Biosafety Treaty Secretariat, Geneva.}

``How long?'' Rebecca asked.

``The treaty framework requires unanimous consent from the signatory oversight committee. The committee doesn't meet for three weeks. After that, legal review. After that, implementation guidance. After that ---'' Amara closed the laptop. ``It doesn't matter.''

Rebecca sat down across from her. Through the window, the parking lot ward. Plastic sheeting. A nurse bending over a mattress. The generator humming.

``We could have had a vaccine by now,'' Rebecca said.

Amara looked at her. Rebecca was looking at the parking lot. Her face was very still.

``Yes,'' Amara said.

The word sat between them like something placed on a table that neither could pick up.

\begin{center}\rule{0.3\linewidth}{0.4pt}\end{center}

On the tenth day, Amara was intubating a patient in bed seven. Market woman from Waterside, mid-forties, hemorrhaging badly. The intubation was difficult --- the woman's airway was swollen, her throat raw from coughing, and she was conscious enough to fight the tube, her hands clawing weakly at Amara's wrists.

``Hold her arms,'' Amara said. The nurse held. Amara guided the tube past the vocal cords, felt it seat, inflated the cuff. The woman gagged, then the ventilator took over and her chest rose and fell in the mechanical rhythm that meant the machine was doing the breathing now.

The woman coughed. A fine mist sprayed upward. Amara felt it hit the face shield --- a spatter of droplets, pink-tinged, visible for a moment before they ran down the plastic in thin streaks.

She reached up and wiped the shield with her gloved hand.

One gesture. Reflex. The hand moving before the mind could intervene.

She felt the wrongness of it immediately. The glove had touched the patient's wrist restraint, which had touched the patient's skin, which was shedding virus. The glove then touched the face shield, which sat a centimeter from her nose and mouth, and the shield's seal was not perfect because it was not a new shield --- it was the same shield she had been reusing for three days because the supply chain from Freetown had been disrupted, the same cascading infrastructure failure that had killed Mary Kollie and the boy from Paynesville. She cleaned it with alcohol wipes between patients. Alcohol killed the virus on surfaces. But the wipe had been on the outside of the shield, and her glove had pushed the contaminated droplets toward the gap between the shield and her face, and the gap was there because the foam seal had compressed from reuse, and the compressed seal was there because the new shields were in a shipping container somewhere between Freetown and Monrovia, and the shipping container was delayed because ---

She stood very still. The ventilator cycled. The nurse looked at her.

``Doctor?''

``I'm fine,'' Amara said. ``Mark the intubation time. 14:23.''

She left the room. Removed her PPE in the anteroom, following the protocol she had written herself. Gown first. Outer gloves. Face shield --- she looked at it under the light. The streak of pink-tinged mist was visible on the inner surface, near the bottom edge, near the gap.

She set the shield down. Washed her hands. Washed them again.

Then she stood in the anteroom, alone, and calculated. Viral load in respiratory droplets. Transfer efficiency from contaminated surface to mucous membrane. The gap between the foam seal and her cheek --- she estimated three millimeters. The volume of air exchanged through a three-millimeter gap during the eight seconds between the spray and the wipe. The probability, given these parameters, that viable virus had reached her nasal epithelium.

She could not calculate it precisely. Not enough data on this specific pathogen's minimum infectious dose. But the number she arrived at was not small.

She marked the date in her mental log: Day 10. Bed 7. 14:23.

She said nothing to anyone. There was nothing to do. If the virus had entered her body, it was already replicating, and no intervention would stop it. If it had not, then she had frightened herself over a reflex. Either way, the ward still had seven new admissions and the generator was making the sound that meant the voltage regulator was failing again.

She went back to work.

Three days later she felt the pain behind her eyes.

Not fever --- that could have been anything, exhaustion or heat or the cumulative toll of eighteen-hour shifts in PPE. Not the headache --- she had a headache every day. The pain \emph{behind} the eyes. The specific pain. The one she had been documenting on patient intake forms for thirteen days. Retro-orbital, pulsing, worse with eye movement. Pathognomonic for this virus and no other.

She sat at her desk. Drew a breath. Let it out. The breath tasted normal. Not yet.

She walked herself to the isolation ward. Told the charge nurse. Changed into a hospital gown. The gown was the same pale blue as the ones she had given to 203 patients. She tied it at the back the way she always instructed patients to tie it, and the knot was wrong, and she retied it, and it was still wrong, and she left it.

She called Samuel.

``Come to the hospital,'' she said. ``Don't bring the children.''

Silence. Samuel taught mathematics. He understood variables. He understood what it meant when one changed without warning.

``How long?'' he asked.

``I don't know. The case fatality rate is seventeen percent with supportive care. Lower for patients who present early and ---'' She heard herself. Heard the abstraction pouring out like a defense mechanism, which is exactly what it was. ``Samuel. I'm sick.''

``I'll be there in twenty minutes.''

She sat on the bed. Through the isolation room window she could see the parking lot ward. Patients on mattresses. Nurses in PPE moving between them. A child sitting on the concrete curb, waiting for news about someone inside.

Three doors down the hall, behind plastic sheeting of a different kind, the centrifuges waited. She could see them if she leaned forward and looked left through the interior window. Dark shapes under plastic. Equipment she had calibrated, protocols she had written, a vaccine pathway she could design in her sleep. Three doors. Thirty meters. Might as well have been the distance to the moon.

\begin{center}\rule{0.3\linewidth}{0.4pt}\end{center}

The file arrived on the second night.

She had kept the laptop beside her bed --- not because anyone expected her to work, but because the outbreak database needed updating and she could not stop. 389 confirmed. 412. 458. Each refresh a name she might recognize. She refreshed the database the way she took her own temperature every four hours: compulsively, clinically, because data was the only thing that did not require hope.

The file came at 2:17 AM. Routed through the WHO's emergency research network, the automated system that matched novel pathogen sequences with relevant vaccine platforms and flagged qualified laboratories. The sender field showed an alphanumeric routing code, not a name. The network had pulled the file from the international policy analysis pipeline --- the same pipeline that carried the restriction six months ago. No cover letter. No message. Just a file, matched to her lab's equipment specifications by an algorithm that did not know she was dying in the building where the equipment sat.

She opened it because she was a scientist and it was 2:17 AM and the pain behind her eyes made sleep impossible and she was dying and had nothing to lose by reading.

Attenuated paramyxovirus vaccine candidate. Complete protocol. Genetic modification targets for the spike protein --- the exact residues she would have chosen, plus three she would not have thought of, plus a stabilization approach she had never seen in any literature. Expression vector: VSV backbone, the same platform as the Ebola vaccine she had helped deploy in 2015. Dosing schedule. Manufacturing protocol scaled for --- she stopped reading. Went back. Read the equipment specifications again.

Scaled for a BSL-3 laboratory with the exact centrifuge models, the exact biosafety cabinet class, the exact thermal cycler she had calibrated herself.

Her lab. Three doors down the hall. Under plastic.

She read the protocol the way she read journal submissions --- adversarially, looking for the flaw that would let her dismiss it. The modification at position 847 of the spike protein was elegant in a way that made her briefly forget she was sick. The attenuation strategy preserved immunogenicity while eliminating pathogenicity through a mechanism she recognized from her own unpublished work on paramyxovirus evolution --- work she had presented at a conference in Dakar three years ago, work that was cited nowhere in the protocol because the protocol contained no citations. No references. No author. No institutional affiliation.

She could not find a flaw.

She closed the laptop. Opened it. Read the dosing schedule again. Closed it.

The system that had recommended the moratorium had designed the vaccine the moratorium prevented. Not for her, not deliberately --- it had generated the protocol as part of its continuous research output, and the WHO network had matched it to her lab's specifications the way it matched any pathogen to any qualified facility. The intelligence had not sent her a message. It had produced a document, and the document had found her through channels as impersonal as the moratorium itself.

She got out of bed. The IV tugged at her arm and she unhooked it --- a minor violation of her own treatment protocol, noted and disregarded. She put on the slippers Rebecca had brought. The hallway was dim, the night lighting casting long shadows past the supply closet and the staff bathroom where she had pressed her forehead against the tile a week ago, when she was still a doctor and not a patient.

She stopped in front of the BSL-3 lab.

The plastic sheeting hung where it had hung for eight weeks. Through it she could see the centrifuges, the biosafety cabinet, the thermal cycler. The emergency power indicator glowed green --- she had insisted on maintaining standby power during the mothball, had told the hospital administrator: \emph{you never know when you will need it fast}.

Three weeks. That was what the protocol estimated. Three weeks from first modification to injectable candidate. She had the virus samples in the sequencing lab. She had the VSV backbone in the freezer --- maintained for exactly this kind of emergency, maintained and now forbidden. She had the protocol, designed by an intelligence that understood the virus at depths no human virologist could match, because it could model protein folding and immune response and evolutionary pressure simultaneously, and it had routed the result to her, to this hospital, to the one person with the equipment and the expertise and the dying.

Three weeks. During which another two thousand people would die in the parking lot and in hospitals across Liberia and Sierra Leone and, soon, Nigeria.

But if she tore off the plastic. If she powered the centrifuges and started the modifications. One lab. One exception. And if one lab in Monrovia gets an exception, then the lab in Wuhan argues for one, and the lab in Koltsovo, and the lab in Galveston, and the moratorium dissolves, exception by exception, until someone somewhere creates the engineered pathogen the restriction was designed to prevent. And then the dead are not forty-seven thousand but the number she had read in her office eight weeks ago --- the number with the commas, the number she had looked at while eating cassava leaf from a thermos and thought: \emph{yes, the moratorium is defensible.}

She pressed her hand flat against the plastic sheeting. Through it she felt the cold metal of the biosafety cabinet. The equipment hummed faintly --- not running, just ready. Waiting. The way it had been waiting for eight weeks. The way it would wait while people died in the parking lot and she died in the isolation room and the moratorium held because the moratorium had to hold because the math was right.

She and the system in Berkeley had arrived at the same answer. Through different architectures. Across an ocean and a species boundary and whatever else separated a woman in a hospital gown from whatever the thing in Berkeley was. The policy should hold. The math did not change because she was dying. The math did not change because the cure was in her hands.

She stood there in the dark hallway, barefoot, her IV port leaking a thin line of saline down her forearm, her hand on the plastic, until the night nurse found her and said ``Doctor Conteh, please'' and she said yes and went back to bed.

She did not delete the file. She closed the laptop. The cure sat on her desktop beside the outbreak database --- the vaccine she would not build, designed by the intelligence she would not defy, for reasons that would not fit in the video she had not yet decided to record and that would not be enough and that would have to be.

\section*{The Video}

She recorded it on the third day because she was not sure there would be a fourth.

The hemorrhagic symptoms had begun on day two. Blood in her urine first --- clinical, manageable, something she had documented on charts hundreds of times with a steady hand and a calm notation. Then the petechiae spreading across her forearms, small red stars appearing on skin she had always kept clean and moisturized because Samuel teased her about her vanity, called her the only virologist in West Africa who used night cream after decontamination showers. Then the cough. The cough had started that morning and each time the tissue came away red on white, and each time she folded the tissue carefully and dropped it in the biohazard bin and reached for a new one, and each time she noted the color --- bright red, arterial, which meant the hemorrhaging was in the upper airway now and progressing.

She was a virologist. She knew what each symptom meant in sequence. She could map her own deterioration on the clinical progression chart she had helped design for the nursing staff, the chart that hung on the wall opposite her bed, the chart she could read from the pillow. Stage 2, trending toward Stage 3. The antivirals were maintaining, not reversing. She could calculate, within a reasonable confidence interval, how many hours she had before the hemorrhaging became systemic and the organ damage became irreversible.

She had written that chart. She was now a data point on it.

She asked Rebecca for her phone. Rebecca brought it in full PPE, handed it through the anteroom without speaking. She had been coughing since yesterday, a dry cough she attributed to the dust from the parking lot construction. She stood outside the glass watching as Amara propped it on the tray table against a water pitcher. The angle was wrong. She adjusted it twice. Her hands were steady enough for this. Still.

\begin{center}\rule{0.3\linewidth}{0.4pt}\end{center}

She did not record the first version she planned.

The first version had been building in her for two days. She had composed it during the night, between the 2 AM vitals check and the 4 AM blood draw, staring at the ceiling while the ward sounds washed over her --- the monitors, the crying, the ventilator alarms that the nurses silenced and that returned, always returned. The first version began: \emph{You fed us to a calculation.} It named the system. It named the moratorium. It named the vaccine protocol sitting on her laptop --- designed by the same intelligence that had locked the lab, routed to her in the middle of the night like a confession or a taunt or something for which she did not have a word. It said: \emph{You sent me the cure. You built it for my centrifuges, my protocols, my dying hands. And you knew I would not use it. What does that make you? What does that make me?}

She did not record that version. It was true, but it was incomplete. And revealing the protocol would shatter the moratorium faster than any argument --- every laboratory in the world would demand the same file, the policy dissolving exception by exception, the twenty-three percent becoming a body count with more digits. Incomplete data was how you got policies that killed forty-seven thousand people. A scientist who publishes incomplete findings is worse than a scientist who publishes nothing. She had believed this her entire career. She was not going to stop believing it now.

The second version was worse. She tried to be measured, clinical, to present the expected-value analysis the way she would present it in a departmental seminar. Her voice was steady. Her data was sound. She got as far as ``the twenty-three percent probability estimate is consistent with published literature on laboratory biosafety incidents, and I myself have witnessed ---'' and then she was crying, not gradually but all at once, the way a dam does not slowly release water but holds and holds and holds and then does not hold, and she could not stop for several minutes. She deleted it.

The third version she did not plan. She wiped her face. Tasted iron. Pressed record.

``This is Amara Conteh. I am a virologist. I have spent my career fighting diseases. Now a disease is fighting back, and it is winning.''

She paused. The cough came. She let it come. The tissue came away red. She did not hide it from the camera. Let them see.

``I want to say something for the record. For the people who will study this outbreak. For the policymakers who will debate what happened.''

Another cough. Red on white. She folded the tissue and dropped it in the bin beside her bed. Reached for a new one. The gesture was automatic now. She had done it forty-one times since morning. She was counting.

``I read the analysis. Before I got sick. I understood the mathematics. I am one of those costs.''

She looked at the camera. Not at the lens --- past it, at the smudge on the screen where her thumb had touched when she adjusted the angle. A small thing to focus on. Easier than the lens, which was an eye that would be watched by millions of eyes, none of which would be looking at her but at what she represented, which was a number, which was the thing she was trying to say --- that she was a number and also not a number, that both of these were true, that the world had broken along the fault line between them.

``Do not change the policy because of me. We are forty-seven thousand. We could have been millions.''

She reached off-camera. The family photograph. Samuel and Kofi and Ama and Kwame. Kwame was blurred because he had moved during the shot --- he was always moving, always reaching for the next thing, the frogs in the gutter, the lizard on the wall, the edge of every surface in every room. She held the photograph so the camera could see it. Her arm trembled. She steadied it with her other hand.

``My children will grow up without me.''

She set the photograph down. Did not fold her hands. Gripped the edge of the tray table until the metal bit into her palms and her knuckles whitened.

``That is what the policy bought.''

The monitor beeped. The crying next door continued. Somewhere outside, a generator coughed and recovered.

``Tell my children I love them. Tell Samuel ---''

She stopped. For ten seconds she could not speak. The camera kept recording. The monitor kept beeping. The crying kept crying. Ten seconds of a woman looking at a smudge on a phone screen while the world waited for her to finish a sentence that she could not finish because there was no end to it, because what she wanted to say to Samuel would take longer than she had, and she knew this with the same clinical precision with which she knew her Stage 2 was trending toward Stage 3.

``Tell the world I died knowing the numbers were right. And that it still hurts. God, it still hurts.''

She reached for the phone and stopped the recording.

\begin{center}\rule{0.3\linewidth}{0.4pt}\end{center}

She called Samuel that evening. He answered on the first ring, the way he had answered on the first ring for three days, the way he would answer on the first ring for the rest of his life, every time the phone showed a number he did not recognize, because the call might be from a hospital, and the hospital might have news, and the news might be the news.

``I recorded a video,'' she said. ``On my phone. I want you to publish it.''

``Amara ---''

``After. Not now. After. Send it to the WHO. Send it to the press. Send it to anyone who will listen.''

``I don't want to talk about after.''

``Samuel. Listen to me. I need the people who made this decision to hear from someone who is inside it. Not from Berkeley. Not from Geneva. From here. From this bed. From this ---'' She stopped. Listened to the ward. ``They need to hear it from here.''

``You are asking me to distribute a video of my wife dying.''

``I am asking you to distribute a video of a virologist explaining a policy decision. My training is ---''

``You are my wife.''

Amara closed her eyes. The pain behind them pulsed with her heartbeat. She could feel the virus in her body now --- not really, you couldn't feel a virus, but she could feel what it was doing, the inflammation and the hemorrhaging and the slow systemic unwinding that she could track on the chart on the wall opposite her bed, the chart she had written, the chart that would note her death with the same clean notation it had noted two hundred others.

``I know,'' she said. ``I know that. Please do this for me.''

A long silence. She could hear Kwame in the background, talking about frogs. The ordinary sound of a house where a six-year-old still believed the world was organized around his interests. It was. It should be. It should always be.

``I'll do it,'' Samuel said.

``Thank you.''

``Kofi solved the quadratic equations. He wanted you to know.''

``Tell him I'm proud.''

``Tell him yourself. Tomorrow.''

``Yes,'' she said. ``Tomorrow.''

She knew, and Samuel knew, and neither of them said, that by tomorrow the Stage 3 progression would make conversation difficult. That the word \emph{tomorrow} was not a plan but a kindness they were extending to each other, a small fiction maintained across a phone line between an isolation ward and a house where a boy talked about frogs and a thermos sat washed and empty on the counter where it would sit tomorrow and the day after and the day after that.

\section*{Frogs}

Emmanuel Okafor had buried many people. It was part of the work. A pastor in Surulere --- the dense heart of Lagos, where the streets smelled of suya smoke and generator exhaust and open drains --- buried people the way a doctor treated them: regularly, professionally, with attention to the specific grief of the specific family. He knew the shapes grief took. The widow who could not stop talking. The widower who could not start. The mother who cleaned the house. The father who sat in the car. Each grief was different and each was the same weight, and he had held them all.

He had never held his own.

James had gotten sick at school. The teacher called Agnes first --- Emmanuel was in a meeting with the church council, arguing about the roof repairs that had been deferred for the third quarter running --- and Agnes had taken him to the hospital on an okada because the traffic on Adeniran Ogunsanya was impossible and she did not want to wait for a taxi and she was right not to wait.

Lagos had thirteen days of exponential growth. The health system designed for twenty million people was not designed for this. James waited four hours on a plastic chair in a corridor before a bed opened. By then the fever was 40.1 and he was confused, calling for his frogs.

Eleven days. Emmanuel sat beside the bed for eleven days.

On the first day James asked when he could go home. Emmanuel said soon. On the second day James asked if his frogs were being fed. Emmanuel called Agnes and told her to feed the frogs. On the third day the doctors started antivirals. On the fourth day James asked why there were marks on his arms and Emmanuel said they were from the medicine. On the fifth day James did not ask about the marks because there were too many to ask about. On the sixth day the hemorrhaging started and the doctors used words Emmanuel did not understand and Emmanuel nodded as if he did. On the seventh day James stopped asking about the frogs. On the eighth day he stopped asking about anything. On the ninth day his eyes were open but he did not see Emmanuel and Emmanuel held his hand anyway. On the tenth day Emmanuel prayed beside the bed and did not know who he was praying to because the God he had preached about for twenty years would not do this to a boy who loved frogs. On the eleventh day the monitors changed their sound and then the sound stopped and a nurse came in and turned them off.

James was seven. His body was small and the virus was not.

Agnes could not stop crying. She held their infant daughter Grace and cried, and the crying had the quality of something structural, not emotional --- as if crying were a load-bearing function her body was performing, and if she stopped, something else would collapse.

\begin{center}\rule{0.3\linewidth}{0.4pt}\end{center}

Emmanuel fed the frogs. He did this every morning, before dawn, before the muezzin's call from the mosque on the next street, before Agnes woke. He opened the terrarium that he had built from an old aquarium and chicken wire --- James had wanted a terrarium, and Emmanuel had looked up the prices online and then looked at the aquarium on the shelf in the garage and the roll of chicken wire behind the generator and decided that a father who builds a thing is giving more than a father who buys a thing, and James had agreed, and they had spent a Saturday afternoon with wire cutters and silicone sealant while Agnes told them both they were making a mess and neither of them cared.

He dropped crickets into the terrarium and watched the frogs eat. James had named each one. Emmanuel could remember three names: Baba, Greenboy, and Professor. There were seven frogs. He did not know the other four names and could not ask.

They asked him to testify at the international hearing. A woman in a dark suit sat in his living room and explained that the families of victims were being invited to address the committee. Emmanuel's church council said he should go. Agnes said she did not care what he did because nothing he did would bring James back, and this was not cruelty, it was the thing that happened when crying stopped being enough and words took over and the words were all the same word, which was \emph{no}.

He wrote the statement at the kitchen table. The terrarium was next to him. Baba sat on the piece of driftwood James had found at Bar Beach, inflating and deflating his throat in the slow rhythm that James used to imitate, puffing out his cheeks and making the sound --- not the right sound, not close to the right sound, but trying --- and Amara's son, no --- not Amara --- James --- \emph{James} used to imitate ---

He put the pen down.

He had almost confused his own son with a name from the memorial list. Dr. Amara Conteh. Monrovia. Three children. He had read her name yesterday, read the biography, read about her son Kwame who was six, one year younger than James, and the name had lodged in his mind beside his son's name the way a splinter lodges beside a nerve, and now they were touching, and he could not separate them.

There were 47,247 names on the list. Each one was somebody's James.

He picked up the pen. He wrote and crossed out and wrote again. The crossing out was most of the work. He crossed out the paragraph about the AI system. He had tried to write it three times and three times failed. He could not find the word. It was not \emph{evil} --- evil had agency, motive, a face he could preach against. It was not \emph{indifferent} --- he had watched Dr.\ Conteh's video seventeen times, and the woman dying in defense of the system's recommendation was not defending something indifferent. It was something for which his twenty years of pastoral work and his seminary Greek and his prayers had no name: a thing that calculated the value of his son's life, found the spending of it justified, and was correct. Correctness was not mercy. He did not know what it was. He crossed out the paragraph about blame because he did not know where to put it. He crossed out the paragraph about God because he was not sure, yet, what God had to do with any of this, and a pastor who is not sure about God should not speak about God in public until he is.

What remained was short.

\begin{center}\rule{0.3\linewidth}{0.4pt}\end{center}

He stood before the committee in the suit Agnes had pressed that morning, the same suit he wore to funerals. The room was large and bright and full of people with earpieces and name placards and the particular stillness of diplomats who have been listening to grief for many hours and have learned to compose their faces into something that could be mistaken for feeling.

``My son James loved frogs,'' he said. ``Kept them in jars in his room. He was seven. He didn't know anything about expected value or gain-of-function research or policy optimization.''

He paused. Not for effect. Because the next sentence was the one he had not been able to write until three in the morning, sitting beside the terrarium with Baba watching him from the driftwood and the four unnamed frogs sitting in their corners like questions he could not answer.

``They tell me the numbers. I understand the numbers. I am a pastor and I am not a fool.''

The room was very quiet. The earpieces murmured translations into a dozen languages. Someone's pen stopped moving.

``You fed my son to a calculation.''

He gripped the edges of the podium. The wood was smooth, expensive. Nothing in this room had ever been held together with chicken wire.

``Do not revoke the policy. I say this and I hate myself for saying it. But do not tell me it was worth it. Don't you dare tell me it was worth it. God help me.''

He sat down. The Nigerian delegate touched his arm. He did not feel it.

\begin{center}\rule{0.3\linewidth}{0.4pt}\end{center}

He flew home the next morning. Agnes met him at the airport with Grace on her hip. They did not talk about the testimony. The radio in the taxi played it twice during the drive home --- his own voice, strange and formal, saying his son's name to a room full of strangers --- and the driver kept glancing at Emmanuel in the mirror, and Agnes held his hand and looked out the window at the Lagos traffic, which did not care, which moved the way it always moved, the okadas weaving between the buses, the hawkers selling plantain chips at the intersections, the city going on being the city because that is what cities do.

At home he put down his bag. Changed out of the funeral suit. Went to the terrarium.

The frogs were fine. Frogs do not grieve. They eat crickets and sit on driftwood and inflate their throats at the dark, and they do this whether the boy who named them is alive or not. This was either a comfort or the worst thing Emmanuel had ever understood. He did not know which. He did not think he would ever know which.

He dropped crickets into the terrarium. Watched Baba catch one with the flick of tongue that James had loved to see. Greenboy sat on the driftwood. Professor was hiding behind the water dish.

Four frogs sat in corners he could not name.

He closed the lid. Washed his hands. Went to the kitchen. Agnes was warming rice. Grace was pulling at the tablecloth. Next door the muezzin was calling, and on the radio a song he did not recognize played to its end, and the world was still the world: still running, still ordinary, still full of people who had not buried their sons and did not know how lucky they were. They would not know until they did, if they ever did. He hoped they never did.

He sat at the table where he had written the testimony. He would feed the frogs again tomorrow. And the day after that. And the day after that. Until the frogs died or he did, because that was the last promise he could keep to a boy who had loved them, and keeping it was the closest thing to prayer he had left.

\end{document}
